\begin{table}[htbp]
    \centering
    \caption{Vergleich der evaluierten Gateway-Hardware}
    \label{tab:gateway_vergleich}
    \begin{tabularx}{\textwidth}{l X X X X}
        \toprule
        \textbf{Gateway}
            & \textbf{Kanäle}
            & \textbf{Formfaktor}
            & \textbf{Stromversorgung}
            & \textbf{Besonderheit} \\
        \midrule
        Dragino LPS8
            & 8
            & Indoor
            & Netzteil
            & PoE-fähig, weit verbreitet \\
        \addlinespace
        MikroTik LR8
            & 8
            & Outdoor
            & PoE / Netzteil
            & Wetterfest (IP54), RouterOS \\
        \addlinespace
        DIY Raspberry Pi
            & 8 (mit RAK2287)
            & Variabel
            & USB / Akku möglich
            & \textbf{Vollständige Kontrolle,}
              \textbf{Network-Server on-board} \\
        \addlinespace
        Pycom LoPy4
            & 1 (single-channel)
            & Kompakt
            & USB / Akku
            & Nur für Tests geeignet \\
        \addlinespace
        TTGO LoRa32
            & 1 (single-channel)
            & Kompakt
            & USB / Akku
            & Nur für Tests geeignet \\
        \bottomrule
    \end{tabularx}
\end{table}

\subsection{Entwicklungs- und Analyse-Tools}
\label{subsec:tools}

Tabelle~\ref{tab:tools} gibt einen Überblick über alle im Rahmen dieser Arbeit
eingesetzten Software"=Werkzeuge.

\begin{table}[htbp]
    \centering
    \caption{Eingesetzte Software-Werkzeuge}
    \label{tab:tools}
    \begin{tabularx}{\textwidth}{l l X}
        \toprule
        \textbf{Bereich} & \textbf{Tool} & \textbf{Zweck} \\
        \midrule
        Firmware-Entwicklung
            & Zephyr RTOS + West
            & RTOS, LoRaWAN"=Stack, Build"=System \\
        \addlinespace
        Simulation
            & Renode
            & Simulation des STM32WL55 ohne Hardware \\
        \addlinespace
        Network"=Server
            & ChirpStack
            & Lokaler, offline"=fähiger LoRaWAN"=Stack \\
        \addlinespace
        Datenpipeline
            & Node"=RED
            & Visuelle Verknüpfung von MQTT"=Daten und Dashboard \\
        \addlinespace
        Debugging
            & MQTT Explorer
            & Inspektion der MQTT"=Nachrichten zwischen ChirpStack und Applikation \\
        \addlinespace
        Spektrumanalyse
            & SDR\# / GNU Radio + \texttt{gr-lora}
            & LoRa"=Chirps visualisieren, Pakete demodulieren \\
        \addlinespace
        Strommessung
            & nRF PPK2 + Power Profiler App
            & Hochauflösende Strom"=Zeit"=Profile \\
        \addlinespace
        RF"=Planung
            & SPLAT!, LoRa Calculator (Semtech)
            & Theoretische Coverage"=Karten, Airtime"=Berechnung \\
        \addlinespace
        Kartenvisualisierung
            & Python + Folium / QGIS
            & RSSI"=Messpunkte als Heatmap auf OpenStreetMap \\
        \addlinespace
        Systemdiagramme
            & draw.io / Lucidchart
            & Blockschaltbilder, Architekturdiagramme \\
        \addlinespace
        Preisvergleich
            & Octopart
            & Komponentenpreise und Verfügbarkeit \\
        \bottomrule
    \end{tabularx}
\end{table}